% Chapter 1: Computing on Opaque Encodings
% DEFINES: ch:opaque-encodings; sec:oblivious-service, sec:untrusted-sees,
%   sec:measured-not-negligible, sec:book-shape, sec:ch1-notes
% REFERENCES (forward, part labels only): part:abstraction,
%   part:confidentiality, part:constructions, part:frontiers

\chapter{Computing on Opaque Encodings}
\label{ch:opaque-encodings}

\section{The Oblivious Service}
\label{sec:oblivious-service}

Imagine a search engine that cannot read your queries, and cannot read the
answers it returns, yet still returns the right ones. You type a word; a
device you trust turns it into an opaque token, a string of bits with no
visible structure. The server holds an enormous lookup table that maps
tokens to tokens. It receives your token, looks up a result token, and
sends it back. Your device turns that result back into an answer. At no
point does the server see the word you searched for, the documents it
returned, or even whether two of your queries were the same. It computes,
usefully, on data whose meaning it cannot recover.

Trace one query through the picture. You search for a sensitive word; call
it $x$. Your device computes a token $c = \enc(x)$, say a sixty-four-bit
string that looks, to anyone without the secret, like a random draw. It
sends $c$ to the server. The server holds a function $\fhat$ that it
applies to $c$ and nothing else: it does not parse $c$, does not look up
$x$ in any plaintext index, does not branch on what $c$ ``means.'' It
returns $\fhat(c)$, another opaque string, and your device decodes it into
the list of matching documents. The server has done real work, a lookup
over a structure that encodes the whole corpus, and has learned nothing it
could write down: not the word, not the result, not even that you searched
twice for the same thing if your device encoded it two different ways.

This is the setting of \emph{trapdoor computing}, and the picture
generalizes well past search. Split the world into two machines. A
\emph{trusted machine} holds a secret and two private operations: an
encoder that turns a value into a token, and a decoder that turns a token
back into a value. We write these $\enc$ and $\dec$, and call the trusted
machine $T$. An \emph{untrusted machine}, which we call $U$, holds only a
function $\fhat$ on bit strings, bits in, bits out, and evaluates it on
whatever it is given. To deploy a latent function $f : X \to Y$ in this
world, the trusted machine encodes an input $x$ as a token $c = \enc(x)$,
the untrusted machine computes $\fhat(c)$ blindly, and the trusted machine
decodes the result as $\dec(\fhat(c))$, an approximation of $f(x)$. The
secret, and the ability to cross between meanings and tokens, lives
entirely on $T$; the untrusted machine sees opaque bits flowing through an
opaque table. Membership testing is the special case $f : X \to \{0,1\}$,
``is this in the set?''; a recommender is $f$ mapping a user token to an
item token; a database lookup is $f$ from a key to a record. The same
two-machine arrangement carries all of them.

Two features of this arrangement do the work, and it is worth fixing them
in mind now because the rest of the book returns to them constantly. The
first is that $\fhat$ is \emph{total}: it is defined on every bit string,
not just on the tokens that encode real inputs. Feed it a random string and
it answers; there is no ``not found,'' no error, no signal that the input
was meaningless. A total function has no special cases for the untrusted
machine to notice. The second is that a single value may have many tokens,
and the tokens are arranged so that, to anyone without the secret, they
look like uniform random strings. Frequency, the usual handle on encrypted
data, finds nothing to grip: the most common query and the rarest one
produce token streams that look alike.

We call $\fhat$, together with its encoder, decoder, and secret, a
\emph{cipher map}: a total function on bit strings that implements a
trapdoor approximation of a hidden function.\footnote{The cipher map is
given its precise definition, as a tuple with four measurable properties,
in \Cref{part:abstraction}. This chapter stays informal on purpose.} The
word ``approximation'' is deliberate and will be paid for honestly: a
cipher map is allowed to be wrong occasionally, and we will see that the
small error budget is not only a construction convenience but a source of
privacy in its own right. The rest of this book is about what such maps can
compute, how their errors compose, how much they leak, and how to build
them. But the whole subject begins with this one image: a machine doing
genuine work on encodings it cannot read.

\section{What the Untrusted Machine Sees}
\label{sec:untrusted-sees}

Put yourself in the position of the untrusted machine. You hold a lookup
table $\fhat$ and you watch a stream of tokens arrive, each of which you
evaluate and return. Over time you see many tokens and many results, and
you may keep every one of them and analyze them at leisure. The governing
question of this book is simple to ask: \emph{how much do you learn?}

The honest answer is that the right design leaves you with very little, and
it is worth being precise about why, because the reasons are not one trick
but four, each blocking a different line of attack.

\begin{itemize}
\item \textbf{It cannot decode.} A token is the image of a value under a
  one-way encoding keyed by a secret the untrusted machine does not hold.
  It sees $c$, never the $x$ behind it, and inverting the encoding is the
  hard direction of a one-way function. This is the most basic wall and the
  one most easily taken for granted.
\item \textbf{It cannot distinguish a real query from filler.} Because
  $\fhat$ is total, the trusted machine can mix in random strings as decoy
  queries, and they produce output exactly as real tokens do. Watch a
  thousand tokens go by; if two hundred of them were decoys, nothing in the
  stream marks which two hundred. There is no ``undefined input'' to give a
  decoy away, and so no way to subtract the noise back out.
\item \textbf{It cannot determine the domain, or which tokens encode the
  same value.} Values are spread across many tokens, and the tokens are
  arranged to be close to uniform, so the frequency analysis that breaks
  deterministic encryption finds no purchase. The untrusted machine cannot
  recover which inputs are in play, and it cannot tell two encodings of one
  value from encodings of two different values; the heavy-hitter query and
  the singleton look the same from where it sits.
\item \textbf{It cannot be sure a result is correct.} A cipher map is
  allowed a small error rate. An occasional wrong answer means that any
  particular answer might be the wrong one, so even a result the untrusted
  machine could somehow read would carry an irreducible ambiguity: it could
  never say ``this membership answer is certainly yes'' rather than ``this
  is one of the cases where the map erred.'' The error budget buys the
  honest user plausible deniability.
\end{itemize}

These four are not accidents of one clever construction. Each is the
consequence of a property that every cipher map has, and that we will name
and measure in \Cref{part:abstraction}: \emph{totality} (every input
produces output, so filler is indistinguishable from real),
\emph{representation uniformity} (tokens are close to uniform, so frequency
analysis fails), \emph{correctness} (a bounded error rate, which doubles as
deniability), and \emph{composability} (the property that lets these maps
be chained, which we will see is also the source of the subtlest leakage).
The pairing is the point: each thing the adversary would like to do is
foreclosed by a named, measurable feature, not by a hope that nobody has
found the attack yet.

What is left to the untrusted machine, then, is exactly a total function on
opaque strings and a log of opaque strings it has evaluated. That is a
remarkably uninformative thing to hold. It is not, however, an
\emph{empty} thing to hold, and the difference between ``uninformative''
and ``empty'' is where the real theory lives. The token stream does have a
distribution; the lookup table does have structure; and a patient observer
can, in principle, estimate things about them. The work of this book is to
say precisely how much, and to give the designer the levers that drive that
amount down. Pretending the leak is zero would be the same mistake the
cryptographic literature has spent two decades correcting in encrypted
search, where ``provably secure'' schemes fell to attacks that read exactly
the residue the proofs declined to measure.

\section{Measured, Not Negligible}
\label{sec:measured-not-negligible}

There is a habit of thought in cryptography that this book does not follow,
and naming it early prevents confusion later. In the usual idiom, a scheme
is secure when an adversary's advantage is \emph{negligible}: any attacker
bounded by reasonable resources succeeds with probability that vanishes
faster than any polynomial in a security parameter, a fact established by
reducing an attack to some problem believed to be hard. Security, in that
idiom, is asymptotic and all-or-nothing. Either the reduction goes through
and the advantage is negligible, or it does not and you have learned
nothing about how bad the leak actually is.

Trapdoor computing asks a different question, and answers it in a different
currency. The untrusted machine here is not trying to ``break'' anything in
the reduction sense; it is a passive observer that gets to watch a
distribution over tokens and is welcome to compute whatever it likes from
it. That distribution is not nothing, and the right response is not to
declare it negligible but to measure it. So instead of asking whether a
bounded adversary can break the scheme, we ask how much the observable
stream actually reveals, in bits, and we hold the designer responsible for
the number. The tools are those of information theory, in the form the
security community calls quantitative information flow: leakage is a
quantity you compute, bound, and then engineer downward.

The central measure is the \emph{entropy ratio} $e = H / H^{*}$: the
entropy of the observed token distribution divided by the most it could
be.\footnote{Here $H^{*} = \log_2 |\mathrm{im}(\enc)|$, the entropy of the
uniform distribution over the tokens the construction actually populates,
and the full theory, including how the parameter $\delta$ that controls
$e$ is reduced and at what cost, is developed in
\Cref{part:confidentiality}.} Read it as a fraction of the way to perfect:
$e = 1$ means the observed stream is as featureless as uniform noise on its
support, and the untrusted machine's best guess about any query is no
better than blind chance; $e < 1$ means the stream carries structure, and
the gap $1 - e$ is the part of an observer's uncertainty that the system
has handed away. The numbers are operational. A system at $e = 0.72$ is
leaking a real fraction of its query structure, enough that an observer
with side information can start to rank guesses; pushed to $e = 0.94$ by
the constructions of \Cref{part:confidentiality}, the residue is small
enough that the same observer is reduced to near-guessing. The designer
moves the number, and pays a stated price in space or bandwidth to do it.

This is what we mean by confidentiality that is \emph{measured} rather than
\emph{negligible}. It is not a stronger claim than the cryptographic one,
nor a weaker one; it is a different and, for this setting, a more useful
one. The cryptographic idiom is the right tool when there is a clean secret
and a clean break to rule out, and it gives a guarantee that holds against
all efficient adversaries at once. The measured idiom is the right tool
when the adversary is handed a total function and asked, honestly, what it
can infer from watching it work; then the actionable answer is a quantity,
not an asymptote, and a quantity is something you can put on a dashboard and
drive toward one.

\section{The Shape of the Book}
\label{sec:book-shape}

Four claims run through everything that follows, and it helps to have them
in one place. First, a single abstraction, the cipher map, \emph{unifies} a
family of constructions that look unrelated at first: encrypted set
membership, oblivious function tables, frequency-hiding retrieval, and the
approximate data structures, Bloom filters and their kin, that the systems
literature already knows. Second, its confidentiality lives at \emph{two
scales} that do not reduce to each other, the leakage of a single map and
the leakage of a chain of maps observed together, and engineering one does
not engineer the other; the second scale, we will find, is created by the
very composability that makes the framework useful. Third, its errors and
its leakage \emph{compose predictably}, which is what makes building real
systems out of these maps tractable rather than hopeful. Fourth, its
privacy comes from \emph{one-way hashing and uniform representation}, not
from hiding access patterns, a distinction sharp enough that the next
chapter is devoted to it.

The book follows those claims in order. \Cref{part:abstraction} builds the
cipher map abstraction: the formal tuple, the four properties, and the
constructions in the abstract. The part after it develops the algebra and
type theory of cipher maps, how Boolean and product and sum structure
survive, or fail to survive, the passage through a trapdoor.
\Cref{part:confidentiality} is the confidentiality theory proper, the two
scales and their measures, and it is where the entropy ratio of the
previous section earns its definition. \Cref{part:constructions} comes down
to earth with concrete constructions, hash-based, linear-algebraic, and
composed, the ones that back the empirical claims made along the way.
\Cref{part:frontiers} closes with the open problems, the places where the
framework runs out and the next program begins. A reader who wants the
whole picture should read it through; a reader who wants one result can
follow the cross-references to the part that holds it, since each is built
to stand on the named properties rather than on the chapter before it.

\section{Notes and Provenance}
\label{sec:ch1-notes}

The framework collected here grew from a set of notes written in 2023 and
2024, which remain the authentic source for its origin: a Bernoulli error
model for approximate computation, an approximate Boolean algebra over
trapdoor sets, an entropy-map construction for hiding function values, and
a treatment of noisy logic gates and how their errors propagate. Those
notes are cited in place as each idea is developed; this book is their
synthesis, and the formal definitions, full proofs, and experiments live in
the later parts and in the standalone papers that those parts cite.

The stance of \Cref{sec:measured-not-negligible}, that confidentiality
should be a measured quantity rather than an asymptotic guarantee, is the
perspective of quantitative information
flow~\cite{smith2009foundations,alvim2020science}, and the
information-theoretic vocabulary used throughout, entropy, conditional
entropy, mutual information, traces back to
Shannon~\cite{shannon1949communication}. A standing convention governs the
whole book: where this synthesis and the formal specification of the
framework disagree on a shared definition, the specification is the
authority, and any divergence found while writing is a bug in the
synthesis, not a new result.
